\documentclass[12pt]{article}

\usepackage{amsmath}
\usepackage{amssymb, amsthm}
\usepackage{geometry}
\usepackage{pifont}
\usepackage{cancel}
\usepackage{newtxtext}
\usepackage[varg]{newtxmath}
\usepackage{enumitem}
\usepackage{setspace}
\usepackage{parskip}
\usepackage{booktabs}
\usepackage{hyperref}
\hypersetup{
    colorlinks=true,
    urlcolor=blue
}

\geometry{a4paper, top=1in, bottom=1in, left=1in, right=1in}

\setstretch{1.15}
\renewcommand{\arraystretch}{1.4}

% 允许页面底部自然留白，不强制拉伸
\raggedbottom

\newenvironment{Answer}{%
  \par\medskip\noindent
  \begin{list}{}{%
    \setlength{\leftmargin}{0pt}%
    \setlength{\itemindent}{0pt}%
    \setlength{\listparindent}{0pt}%
  }%
  \item\relax
  \textbf{Answer.}\quad\ignorespaces
}{%
  \end{list}
  \par\medskip
}

\setlist[enumerate]{itemsep=0.8em, topsep=0.8em, leftmargin=*}
\setlist[itemize]{itemsep=0.4em, topsep=0.4em, leftmargin=*}

\begin{document}

\begin{center}
    {\Large\bfseries SHENZHEN UNIVERSITY} \\[0.3em]
    {\large\bfseries INSTITUTE FOR ADVANCED STUDY} \\[0.8em]
    \textit{Differential Equations, Autumn 2026} \\
    Assignment 1, due at 1:00 PM on Wednesday, September 30, 2026
\end{center}

\vspace{1.5em}

\section*{1.1}
In Problem 13-15, write a differential equation that fits the physical description.
\begin{enumerate}[label=\textbf{\arabic*.}]
  \item[13.] The rate of change of the population $p$ of bacteria at time $t$ is proportional to the population at time $t$.
  \begin{Answer}
  \begin{equation*}
    \displaystyle \frac{dp}{dt} = kp
  \end{equation*}
  \end{Answer}

  \item[14.] The velocity at time $t$ of a particle moving along a straight line is proportional to the fourth power of its position $x$.
  \begin{Answer}
  \begin{equation*}
    v = kx^4
  \end{equation*}
  \end{Answer}
\end{enumerate}

\section*{1.2}
\begin{enumerate}[label=\textbf{\arabic*.(a)}, leftmargin=*]
  \item Show that $\phi(x)=x^{2}$ is an explicit solution to
  \[ \displaystyle x\frac{dy}{dx}=2y \]
  on the interval $(-\infty,\infty)$.
  \begin{Answer}
  \begin{align*}
    \phi(x) &= x^2, \quad x \in (-\infty, \infty) \\[0.3em]
    \displaystyle x \cdot \frac{dy}{dx} &= \displaystyle x \cdot (x^2)' = x \cdot 2x = 2 \cdot x^2 = 2\phi(x) = 2y
  \end{align*}
  \end{Answer}

  \item[(b)] Show that $\phi(x)=e^{x}-x$ is an explicit solution to
  \[ \displaystyle \frac{dy}{dx}+y^{2}=e^{2x}+(1-2x)e^{x}+x^{2}-1 \]
  on the interval $(-\infty,\infty)$.
  \begin{Answer}
  \begin{align*}
    \phi(x) &= e^x - x, \quad x \in (-\infty, \infty) \\[0.3em]
    \displaystyle \frac{dy}{dx} + y^2 &= (e^x - x)' + (e^x - x)^2 \\[0.3em]
    &= e^x - 1 + e^{2x} - 2e^x x + x^2 \\[0.3em]
    &= e^{2x} + (1-2x)e^x + x^2 - 1
  \end{align*}
  \end{Answer}

  \item[(c)] Show that $\phi(x)=x^{2}-x^{-1}$ is an explicit solution to $\displaystyle x^{2}\frac{d^{2}y}{dx^{2}}=2y$ on the interval $(0,\infty)$.
  \begin{Answer}
  \begin{align*}
    \phi(x) &= x^2 - x^{-1}, \quad x \in (0, \infty) \\[0.3em]
    \displaystyle x^2 \frac{d^2y}{dx^2} &= \displaystyle x^2 \cdot (x^2 - x^{-1})'' = x^2 \cdot (2x + x^{-2})' \\[0.3em]
    &= x^2 \cdot (2 - 2x^{-3}) \\[0.3em]
    &= 2(x^2 - x^{-1}) \\[0.3em]
    &= 2\phi(x) = 2y
  \end{align*}
  \end{Answer}
\end{enumerate}

\section*{1.2 (Continued)}
\begin{enumerate}[label=\textbf{\arabic*.}, leftmargin=*]
  \item[3.] $\displaystyle y=\sin x+x^{2},\ \frac{d^{2}y}{dx^{2}}+y=x^{2}+2$.
  \begin{Answer}
  \begin{align*}
    y &= \sin x + x^2 \\[0.3em]
    y' &= \cos x + 2x \\[0.3em]
    y'' &= -\sin x + 2 \\[0.3em]
    \displaystyle \frac{d^2y}{dx^2} + y &= -\sin x + 2 + \sin x + x^2 = x^2 + 2
  \end{align*}
  \end{Answer}

  \item[8.] $\displaystyle y=3\sin 2x+e^{-x},\ y^{\prime\prime}+4y=5e^{-x}$.
  \begin{Answer}
  \begin{align*}
    y &= 3\sin 2x + e^{-x} \\[0.3em]
    y' &= 6\cos 2x - e^{-x} \\[0.3em]
    y'' &= -12\sin 2x + e^{-x} \\[0.3em]
    y'' + 4y &= -12\sin 2x + e^{-x} + 4(3\sin 2x + e^{-x}) \\[0.3em]
    &= -12\sin 2x + e^{-x} + 12\sin 2x + 4e^{-x} \\[0.3em]
    &= 5e^{-x}
  \end{align*}
  \end{Answer}

  \item[10.] $\displaystyle y-\ln y=x^{2}+1,\ \frac{dy}{dx}=\frac{2xy}{y-1}$.
  \begin{Answer}
  \begin{align*}
    y - \ln y &= x^2 + 1 \\[0.3em]
    \displaystyle dy - \frac{1}{y}dy &= 2x \, dx \\[0.3em]
    \displaystyle \frac{dy}{dx} - \frac{1}{y}\frac{dy}{dx} &= 2x \\[0.3em]
    \displaystyle \left(1 - \frac{1}{y}\right)\frac{dy}{dx} &= 2x \\[0.3em]
    \displaystyle \frac{dy}{dx} &= \displaystyle \frac{2x}{1 - \frac{1}{y}} = \frac{2x}{\frac{y-1}{y}} = \frac{2xy}{y-1}
  \end{align*}
  \end{Answer}

  \item[11.] $\displaystyle e^{xy}+y=x-1,\ \frac{dy}{dx}=\frac{e^{-xy}-y}{e^{-xy}+x}$.
  \begin{Answer}
  \begin{align*}
    e^{xy} + y &= x - 1 \\[0.3em]
    d(e^{xy}) + dy &= dx \\[0.3em]
    e^{xy}(x \, dy + y \, dx) + dy &= dx \\[0.3em]
    \displaystyle e^{xy} \cdot x \cdot \frac{dy}{dx} + e^{xy} \cdot y + \frac{dy}{dx} &= 1 \\[0.3em]
    \displaystyle \frac{dy}{dx}(1 + e^{xy} \cdot x) &= 1 - e^{xy} \cdot y \\[0.3em]
    \displaystyle \frac{dy}{dx} &= \displaystyle \frac{1 - e^{xy} \cdot y}{1 + e^{xy} \cdot x} = \frac{e^{xy}(e^{-xy} - y)}{e^{xy}(e^{-xy} + x)} \\[0.3em]
    &= \displaystyle \frac{e^{-xy} - y}{e^{-xy} + x}
  \end{align*}
  \end{Answer}

  \item[20.] Determine for which values of $m$ the function $\phi=e^{mx}$ is a solution to the given equation.
  \[ \begin{array}{l}
    (a)\ \displaystyle \frac{d^{2}y}{dx^{2}}+6\frac{dy}{dx}+5y=0, \\[4pt]
    (b)\ \displaystyle \frac{d^{3}y}{dx^{3}}+3\frac{d^{2}y}{dx^{2}}+2\frac{dy}{dx}=0.
  \end{array} \]
  \begin{Answer}
  \textbf{(a)} \quad \(\phi(x) = e^{mx}\)
  \begin{align*}
    \displaystyle \frac{dy}{dx} &= m e^{mx}, \quad \frac{d^2y}{dx^2} = m^2 e^{mx} \\[0.4em]
    \displaystyle \frac{d^2y}{dx^2} + 6\frac{dy}{dx} + 5y &= m^2 e^{mx} + 6m e^{mx} + 5e^{mx} \\[0.3em]
    &= m^2 e^{mx} + (6m + 5)e^{mx} = 0 \\[0.4em]
    m^2 + 6m + 5 &= 0 \\[0.3em]
    m &= \frac{-6 \pm \sqrt{36 - 4}}{2} = -3 \pm 2\sqrt{2}
  \end{align*}

  \vspace{0.2em}
  \textbf{(b)} \quad \(\displaystyle \frac{d^3y}{dx^3} = m^3 e^{mx}\)
  \begin{align*}
    \displaystyle \frac{d^3y}{dx^3} + 3\frac{d^2y}{dx^2} + 2\frac{dy}{dx} &= 0 \\[0.3em]
    m^3 e^{mx} + 3m^2 e^{mx} + 2m e^{mx} &= 0 \\[0.3em]
    m^3 + 3m^2 + 2m &= 0
  \end{align*}
  \ding{172} If \(m = 0\): \quad \(m^3 + 3m^2 + 2m = 0^3 + 3 \cdot 0^2 + 2 \cdot 0 = 0\)
  
  \(m = 0\) is a solution. \par
  \vspace{0.1em}
  \ding{173} If \(m \neq 0\): \quad
  \begin{align*}
    m^2 + 3m + 2 &= 0 \\[0.3em]
    (m+2)(m+1) &= 0 \\[0.3em]
    m &= -2 \text{ or } -1
  \end{align*}
  Thus, \(m = 0 \text{ or } -2 \text{ or } -1\).
  \end{Answer}

  \item[22.] Verify that the function $\phi(x)=c_{1}e^{x}+c_{2}e^{-2x}$ is a solution to the linear equation
  \[ \displaystyle \frac{d^{2}y}{dx^{2}}+\frac{dy}{dx}-2y=0 \]
  for any choice of the constants $c_{1}$ and $c_{2}$. Determine $c_{1}$ and $c_{2}$ so that each of the following initial conditions are satisfied.
  \begin{itemize}[leftmargin=*]
      \item[(a)] $y(0)=2,\ y^{\prime}(0)=1$;
      \item[(b)] $y(1)=1,\ y^{\prime}(1)=0$.
  \end{itemize}
  \begin{Answer}
  \begin{align*}
    \phi(x) &= c_1 e^x + c_2 e^{-2x} \\[0.3em]
    \displaystyle \frac{dy}{dx} &= c_1 e^x - 2c_2 e^{-2x} \\[0.3em]
    \displaystyle \frac{d^2y}{dx^2} &= c_1 e^x + 4c_2 e^{-2x} \\[0.4em]
    \displaystyle \frac{d^2y}{dx^2} + \frac{dy}{dx} - 2y &= c_1 e^x + 4c_2 e^{-2x} + c_1 e^x - 2c_2 e^{-2x} - 2(c_1 e^x + c_2 e^{-2x}) \\[0.3em]
    &= (c_1 + c_1 - 2c_1)e^x + (4c_2 - 2c_2 - 2c_2)e^{-2x} = 0
  \end{align*}
  So \(\phi(x) = c_1 e^x + c_2 e^{-2x}\) is a solution, \(c_1, c_2 \in \mathbb{R}\).

  \vspace{0.3em}
  \textbf{(a)} \quad
  \begin{align*}
    y(0) &= c_1 e^0 + c_2 e^0 = c_1 + c_2 = 2 \\[0.3em]
    y'(0) &= c_1 e^0 - 2c_2 e^0 = c_1 - 2c_2 = 1
  \end{align*}
  \[
  \left( \begin{array}{cc|c}
    1 & 1 & 2 \\
    1 & -2 & 1
  \end{array} \right)
  \xrightarrow{\text{\ding{173}} + \text{\ding{172}} \cdot (-1)}
  \left( \begin{array}{cc|c}
    1 & 1 & 2 \\
    0 & -3 & -1
  \end{array} \right)
  \quad \left\{ \begin{aligned}
    c_1 &= \frac{5}{3} \\
    c_2 &= \frac{1}{3}
  \end{aligned} \right.
  \]

  \vspace{0.5em}
  \textbf{(b)} \quad
  \begin{align*}
    y(1) &= c_1 e^1 + c_2 e^{-2} = 1 \\[0.3em]
    y'(1) &= c_1 e^1 - 2c_2 e^{-2} = 0
  \end{align*}
  \[
  \left( \begin{array}{cc|c}
    e & e^{-2} & 1 \\
    e & -2e^{-2} & 0
  \end{array} \right)
  \xrightarrow{\text{\ding{173}} + \text{\ding{172}} \cdot (-1)}
  \left( \begin{array}{cc|c}
    e & e^{-2} & 1 \\
    0 & -3e^{-2} & -1
  \end{array} \right)
  \quad \left\{ \begin{aligned}
    c_1 &= \frac{2}{3e} \\
    c_2 &= \frac{1}{3}e^2
  \end{aligned} \right.
  \]
  \end{Answer}
\end{enumerate}

\section*{6.}
Use Euler's method with step size $h=0.2$ to approximate the solution to the initial value problem
\[ \displaystyle y^{\prime}=\frac{1}{x}\left(y^{2}+y\right),\quad y(1)=1 \]
at the points $x=1.2, 1.4, 1.6$, and $1.8$.
\begin{Answer}
\begin{align*}
  y' &= f(x, y) = \displaystyle \frac{1}{x}(y^2 + y) \\[0.3em]
  y_{n+1} &= y_n + h \cdot f(x_n, y_n) \\[0.5em]
  y(1.2) &= y(1) + h \cdot f(1, 1) = 1 + 0.2 \times \left[ 1 \times (1^2 + 1) \right] \\[0.3em]
  &= 1 + 0.2 \times 2 = 1.4
\end{align*}

\vspace{0.5em}
\begin{align*}
  y(1.4) &= y(1.2) + h f(1.2, 1.4) \\[0.3em]
  &= 1.4 + 0.2 \times \left[ \displaystyle \frac{1}{1.2} \times (1.4^2 + 1.4) \right] \\[0.3em]
  &= 1.4 + 0.2 \times 2.8 = 1.96
\end{align*}

\vspace{0.5em}
\begin{align*}
  y(1.6) &= y(1.4) + h f(1.4, 1.96) \\[0.3em]
  &= 1.96 + 0.2 \times \left[ \displaystyle \frac{1}{1.4} \times (1.96^2 + 1.96) \right] \\[0.3em]
  &= 1.96 + 0.2 \times 4.144 = 2.7888
\end{align*}

\vspace{0.5em}
\begin{align*}
  y(1.8) &= y(1.6) + h f(1.6, 2.7888) \\[0.3em]
  &= 2.7888 + 0.2 \times \left[ \displaystyle \frac{1}{1.6} \times (2.7888^2 + 2.7888) \right] \\[0.3em]
  &= 2.7888 + 0.2 \times 6.6039 = 4.1096
\end{align*}
\end{Answer}

\vspace{0.8em}
\begin{center}
\begin{tabular}{c c}
\toprule
\(x\) & \(y\) \\
\midrule
1.2 & 1.4 \\
1.4 & 1.96 \\
1.6 & 2.7888 \\
1.8 & 4.1096 \\
\bottomrule
\end{tabular}
\end{center}

\vspace{0.5em}
\noindent
\href{https://github.com/BrianCai666/Differential-Equations-2026Fall/blob/main/hw/hw1/euler.c}{Code implementation: euler.c}
\end{document}